\cotefile{81}
\scanpage{35}
d'où
\[
\Sigma_{12}=X\deleted{\unsure{-}}\unsure{\setminus|-}(\Sigma_1^+\cup\Sigma_2^+)
\]
\unsure{donc} $\Sigma_{12}=\varnothing$ signifierait que
\[
X=\Sigma_1^+\amalg\Sigma_2^+
\]
et par suite
\[
\Pi_1=\left\{\Sigma_1^+,\underset{\inserted{\substack{\Vert\\\Sigma_{\unsure{1|2}}^-=X\setminus\Sigma_1^+}}}{\Sigma_2^+}\right\},
\quad \Pi_2=\{\Sigma_2^+,\Sigma_1^+\}
\]
alors $\deleted{\unsure{\Pi_2|\Pi_1}}\quad\Pi_1=\Pi_{\unsure{3|2}}$ contrairement \unsure{:|à} l'\unsure{hypothèse}.

On trouve donc

\noindent\makebox[0pt][r]{\rule[-2.7ex]{.4pt}{3.8ex}\hspace{.45em}}\deleted{\uline{Corollaire}} \inserted{Proposition} \unsure{Soient|Soit} $\Pi_1\unsure{;|,}\Pi_2$ deux décompositions binaires de $X$, mutuellement non \unsure{panachées}.
% Vertical emphasis strokes to the left of “binaires … Alors” and “tel”.
Alors il existe un \uline{unique} couple
\[
(\Sigma_1^{\unsure{+|\top}},\Sigma_2^+)\in\Pi_1\times\Pi_2
\]
\makebox[0pt][r]{\rule[-.35ex]{.4pt}{2ex}\hspace{.45em}}tel que l'on ait
\[
\Sigma_1^+\cap\Sigma_2^+=\varnothing
\]

\begingroup\makeatletter\tagsleft@true\makeatother
Soit alors
\[
C=\Pi_1\amalg\Pi_2=\unsure{\omega_1|w_1}\amalg\omega_2
\tag*{(17)}
\]
l'ensemble somme, qui a quatre éléments. On a une application canonique
\[
C\longrightarrow{\unsure{\mathfrak{P}|\mathcal{P}}}_{\unsure{\mathrm{finies}|\mathrm{ferm\acute{e}s}|\mathrm{foncier}}}(X)
\tag*{(1\deleted{\unsure{7}}8)}
\]
qui a comme image
\[
\{\Sigma_1^+,\Sigma_1^-,\Sigma_2^+,\Sigma_2^-\},
\tag*{(1\deleted{\unsure{8}}9)}
\]
\unsure{et|à} laquelle on l'identifie si $(1\deleted{\unsure{7}}8)$ est \unsure{injective}, i.e.

\endgroup
\clearpage
% archive sheet number: 34
